\subsubsection{圆方树 (Round-Square Tree)}
圆方树 (Block-Cut Tree)，用于处理一般图上的点双连通分量。
\begin{lstlisting}
ll n,m,tim;
ll dfn[100010],low[100010];
ll st[100010],tp,c;
bool in[100010];
vi e[100010];
vi E[200010];
void dfs(ll u,ll fa){
	dfn[u]=low[u]=++tim;
	st[++tp]=u;in[u]=1;
	bool skipped=0;
	fe(i,e[u]){
		ll v=e[u][i];
        if(v==fa&&!skipped){skipped=1;continue;}
		if(!dfn[v]){
			dfs(v,u);low[u]=min(low[u],low[v]);
			if(low[v]>=dfn[u]){
				c++;E[c].pb(u),E[u].pb(c);
				ll now;while(1){
					now=st[tp--];in[now]=0;
					E[now].pb(c),E[c].pb(now);
					if(now==v)break;
				}
			}
		}else if(in[v])low[u]=min(low[u],dfn[v]);
	}
}

int main(){
	n=c=d,m=d;tim=tp=0;
    f(i,1,n)e[i].clear(),dfn[i]=low[i]=in[i]=0;
    f(i,1,2*n)E[i].clear();
	f(i,1,m){ll u=d,v=d;e[u].pb(v),e[v].pb(u);}
	f(i,1,n)if(!dfn[i]){dfs(i,0);tp=0;in[i]=0;}
	return 0;
}
\end{lstlisting}

\paragraph{多测与注意} 保留原顶点栈写法，已补多测清空并只跳一条父边。非连通图得到森林；方点从 n+1 编号。
